% Source: https://artofproblemsolving.com/community/c3h1606763_i_know_how_to_do_this_but_i_just_have_a_question
% Stable identity: aops-topic-1606763-q01
By the triangle inequality (strict, because the area is positive), a triangle with sides $4$, $6$, $x$ exists exactly when $6-4<x<6+4$, i.e.\ $2<x<10$; likewise $2<y<10$.

\begin{asy}
usepackage("amsmath");
size(8cm);
defaultpen(fontsize(10pt));
pair V=(0,0);
pair B1=(6,0),A1=4*dir(25);
draw(V--B1--A1--cycle,linewidth(0.9));
label("$6$",(V+B1)/2,S);label("$4$",(V+A1)/2,NW);label("$x$",(A1+B1)/2,NE,blue);
pair W2=(0,-4);
pair B2=W2+(6,0),A2=W2+4*dir(155);
draw(W2--B2--A2--cycle,linewidth(0.9));
label("$6$",(W2+B2)/2,S);label("$4$",(W2+A2)/2,SW);label("$y$",(A2+B2)/2,N,blue);
label("\mbox{}",(min(currentpicture,user=true)+max(currentpicture,user=true))/2,black); // keeps the text after the figure black when LaTeX typesets the labels inline
\end{asy}

\textbf{Which differences occur.} If $2<x,y<10$, then $|x-y|<10-2=8$. Conversely every value $0\le t<8$ occurs: take $x=2+\varepsilon$ and $y=2+\varepsilon+t$ with $0<\varepsilon<8-t$, so both lie in $(2,10)$.

So the possible values of $|x-y|$ are exactly $[0,8)$, and the smallest positive number that is not possible is $8$:
\[
\boxed{\textbf{(D)}\ 8}.
\]
(The figure shows a narrow triangle, with $x$ close to $2$, and a wide one, with $y$ close to $10$; their difference is close to $8$ but can never reach it.)
